\section{Introduction}

A Hopfield network stores patterns as the stable states of a recurrent network and recalls one from a corrupted copy \cite{hopfield1982}. How many patterns it can hold is the classic question. For independent random patterns and Hebbian weights, the answer is a load of roughly $0.138N$ patterns for $N$ neurons if a small fraction of wrong bits is tolerated \cite{amit1985}, and fewer if every bit must come back \cite{mceliece1987}. Real patterns are rarely independent. In the assignment's character bitmaps, most pixels are on, and the bitmaps share strokes, so stored patterns overlap, and overlap is the hard case.

The assignment stored \N{d:n_patterns} dense $12 \times 12$ character bitmaps in a network of \N{d:n_neurons} neurons with Hebbian weights and tried to recall them from copies with about 10\% and 20\% of the pixels flipped. The baseline recalled none of them. This paper reports that baseline, corrected, and asks three questions of it. How fast does exact recall collapse as stored patterns grow alike? Which learning rules hold up? Why do the characters fail?

The paper makes three contributions: the corrected baseline (Section~\ref{sec:baseline}); a capacity study on random patterns with controlled overlap and density, which compares nine recall rules (Sections~\ref{sec:design} and~\ref{sec:results}); and the placement of the characters on the synthetic curves (Section~\ref{sec:why}).

\emph{Revision note.} The original coursework submission (spring 2026) stored the patterns in an unsigned 8-bit type, so turning a 0 into $-1$ wrapped around to 255 and the summed weights came out wrong (from 0 to 4,082). Its baseline numbers (every query ending in the all-ones state, a final distance of 15 bits to the nearest stored pattern, and the energies) came from that error. This revised version (fall 2026) replaces them with a rerun with correct weights, and every single-run number from the original improvements is replaced by a distribution or an interval. The energy now follows the assignment's definition on 0/1 states, and the capacity study is new. Other methods from the original submission are left out, each for its own reason. The threshold and mean-centering corrections recalled 0 of 18 everywhere. Four of the five density normalizations left the patterns unchanged, and the fifth, target-density mapping, lowered the pseudoinverse rule to 3 of 18 on the test file. The sparse codes were scored by nearest pattern rather than exact recall, so their numbers do not compare. The subset selection was a single run (at best 1 of 3 patterns recalled), and the random subsets of Section~\ref{sec:why} replace it.
